\section{Benchmarks}\label{sec:benchmarks}

\subsection{Research Questions}
\guide{What the experiments answer. E.g.\ RQ1: do the methods agree and verify? RQ2: how much does
knowledge shrink the product? RQ3: speed vs.\ \tool{ltlfsynt}? RQ4: which method wins where?}

\begin{enumerate}
  \item[RQ1] \guide{question}
  \item[RQ2] \guide{question}
\end{enumerate}

\subsection{Metrics}
\guide{Structural (state counts, deterministic, committed) vs.\ timing (wall clock, median of
repeats). Prune ratio; nk-product ratio; speedup. Why state counts are the decisive number.}

\subsection{Baselines}
\guide{\tool{ltlfsynt} on the $\psiin\rightarrow\varphi$ encoding; the no-knowledge product;
the unpruned product. Comparability tiers: which families admit a fair external comparison.}

\subsection{Benchmark Families}
\guide{One paragraph per family: $\varphi_n$, knowledge, tier, what it isolates. Then the
summary table.}

\subsubsection{Parametric Families}
\guide{cons-prunes, cons-inert, mirror-small, mirror-degenerate, parity-t3, knowledge-chain(-inert),
co-generated, \ldots}

\subsubsection{Engineered Domain Families}
\guide{Domain-inspired instances (e.g.\ slippery world, hidden counters).}

\subsubsection{External Suites}
\guide{SYNTCOMP finite-trace instances; what fraction carries usable knowledge.}

\subsubsection{Generated Corpus}
\guide{Random instances used for oracle soak testing.}

\begin{table}[H]
\centering
\caption{Benchmark families.}\label{tab:families}
\begin{tabular}{lllll}
\toprule
Family & $\varphi_n$ & Knowledge & Tier & Purpose \\
\midrule
 & & & & \\
\bottomrule
\end{tabular}
\end{table}

\subsection{Experimental Setup}
\guide{Machine, build type, Spot version, commit, repeats, timeout, budget.}
